Los requisitos funcionales del \dfn{frontend} se obtuvieron aplicando el procedimiento de
\cref{SEC:AN-BACKEND} a las operaciones del contrato. No se parte de una lista de deseos:
cada requisito corresponde a operaciones existentes del \dfn{backend}, y lo que la interfaz
necesitaría pero el contrato no permite se trata aparte como hueco (\cref{SEC:AN-GAPS}). Los
requisitos del servicio ya tienen su propia numeración (RF-1 a RF-16 en la memoria y en los
resúmenes del OpenAPI); los del \dfn{frontend} se numeran como RF-F para no confundirlos, y
cada uno indica las etiquetas del contrato que usa. El detalle operación por operación está en
el \cref{AP:API}.

\Cref{TAB:AN-RF} resume los requisitos, agrupados por el actor principal y en el orden del
ciclo de vida de un examen. La última columna indica en qué línea de trabajo se implementó:
\emph{D} para los flujos docentes y del estudiante, \emph{A} para la línea paralela de
administración, cuenta y notificaciones (\cref{SEC:IMP-PARALLEL}).

\begin{table}[htbp]
  \centering
  \small
  \caption{Requisitos funcionales del \textit{frontend} y su trazabilidad al contrato.}
  \label{TAB:AN-RF}
  \begin{tabular}{p{1.3cm}p{6.6cm}p{3.6cm}c}
    \hline
    \textbf{Id.} & \textbf{Requisito} & \textbf{Etiquetas del contrato} & \textbf{Línea} \\
    \hline
    \multicolumn{4}{l}{\textit{Todos los actores}} \\
    RF-F1  & Iniciar y cerrar sesión; renovar el \acs{jwt} de forma transparente y volver
             al inicio de sesión si la renovación falla. & Authentication & base \\
    RF-F2  & Ver un inicio con las tareas propias (correcciones y revisiones pendientes) y
             las asignaturas con el rol en cada una. & Grading, Subjects & D \\
    RF-F3  & Consultar y editar el perfil propio y sus preferencias de notificación;
             consultar y marcar notificaciones. & Users, Notifications & A \\
    \multicolumn{4}{l}{\textit{Coordinación}} \\
    RF-F4  & Crear y configurar exámenes: modelos con su \acs{pdf} en blanco, problemas,
             permisos del estudiante y zonas de reconocimiento por página. & Exams & D \\
    RF-F5  & Definir \dfnpl{rubrica} y reglas de asignación de correctores. & Grading,
             Exams & D \\
    RF-F6  & Resolver incidencias de escaneo: identificar estudiante o modelo, aceptar,
             reordenar, mover o descartar páginas, adjuntar páginas huérfanas, crear o
             borrar \dfnpl{instancia}. & Instances, Pages & D \\
    RF-F7  & Seguir el estado de las \dfnpl{instancia} de un examen, publicar notas y
             exportarlas. & Instances, Grading, Exams & D \\
    RF-F8  & Abrir el periodo de revisión y gestionar sus solicitudes. & Reviews & D \\
    RF-F9  & Administrar miembros, grupos y permisos efectivos de la asignatura.
           & Subjects & A \\
    \multicolumn{4}{l}{\textit{Docente}} \\
    RF-F10 & Corregir una \dfn{instancia}: ver el examen escaneado, anotar a mano alzada,
             calificar cada problema manualmente o por \dfn{rubrica}, comentar y cambiar
             su estado. & Instances, Annotations, Grading & D \\
    RF-F11 & Resolver las solicitudes de revisión asignadas. & Reviews & D \\
    \multicolumn{4}{l}{\textit{Estudiante}} \\
    RF-F12 & Consultar la nota publicada y, según los permisos del examen y el periodo de
             revisión, el desglose, la \dfn{rubrica} y el examen anotado. & Instances,
             Annotations & D \\
    RF-F13 & Solicitar la revisión de problemas concretos y seguir su resolución.
           & Reviews & D \\
    \multicolumn{4}{l}{\textit{Gestor de la organización}} \\
    RF-F14 & Administrar usuarios (alta, edición, importación y exportación, restablecer
             contraseña), cursos académicos, asignaturas, configuración y auditoría.
           & Users, Courses, Organizations, Subjects, Audit & A \\
    RF-F15 & Ingesta manual de escaneos y atención de páginas sin examen asociado.
           & Ingestion, Pages & D \\
    \hline
  \end{tabular}
\end{table}

\subsection{Requisitos transversales de comportamiento}

Además de las funciones, el análisis del servicio impone reglas de comportamiento que afectan
a varias pantallas a la vez:
\begin{description}
  \item[RF-F16. Variantes por rol y permiso.] Cada vista se adapta al rol de la persona en la
        asignatura y a sus permisos: un examen es una pantalla de seguimiento para la
        coordinación, de trabajo para el docente y de resultado para el estudiante.
  \item[RF-F17. Estados de carga.] Toda lista o detalle remoto distingue los estados de
        carga, vacío, error y falta de permiso, y permite reintentar.
  \item[RF-F18. Conflictos de edición.] Al recalificar, la interfaz envía la nota anterior
        conocida y, ante un conflicto de versión, avisa en lugar de sobrescribir.
  \item[RF-F19. Persistencia de anotaciones.] Los trazos se guardan en el \dfn{backend} sin
        acción explícita y no se pierden al cambiar de examen o salir del corrector.
  \item[RF-F20. Errores comprensibles.] Los códigos de error simbólicos del servicio se
        traducen a mensajes en el idioma de la interfaz.
\end{description}

\subsection{Alcance respecto al contrato}

Con las dos líneas de trabajo reunidas, el código de las funcionalidades invoca 85 de las 94
operaciones funcionales del contrato, y el cliente común realiza otras dos (renovación y
cierre de sesión), según el \textit{script} de cobertura ejecutado sobre ambas ramas el
18/10/2026. Quedan sin interfaz siete operaciones, por motivos distintos:
\begin{itemize}
  \item La creación de organizaciones, reservada al superadministrador, que queda fuera de
        los actores de la interfaz.
  \item Crear, modificar y borrar asignaturas mediante sus operaciones genéricas, que la
        administración de la organización no ha llegado a usar.
  \item Las dos operaciones de reconocimiento de texto sobre anotaciones (\acs{ocr} de trazos
        y de notas manuscritas), previstas como ayuda al corrector.
  \item El adjunto de páginas huérfanas por su ruta propia, que falla siempre en la versión
        analizada del servicio; la interfaz usa una operación equivalente que sí funciona
        (\cref{SEC:AN-GAPS}).
\end{itemize}
La medición, su método y su evolución se presentan en \cref{SEC:TEST-COVERAGE}.
